\begin{proof}[Proof of \cref{obj:prop:signed-causal-bridge}]
Throughout, \(d\ge2\), so every stratum has positive probability under a law in \(\mathcal V_d\).

\begin{enumerate}
\item \emph{Identification of the signed mean.}
Fix \(P\in\mathcal V_d\) and \(j\in[d]\). By \cref{obj:ass:fair-randomization}, summing over the potential outcome for the other arm gives, for \(a\in\{0,1\}\),
\[
\begin{aligned}
&P(X=j,A=a,Y^a=1)\\
&\qquad=\sum_{y\in\{0,1\}}
 P(X=j,A=a,Y^a=1,Y^{1-a}=y)\\
&\qquad=\frac12\sum_{y\in\{0,1\}}
 P(X=j,Y^a=1,Y^{1-a}=y)
 =\frac12P(X=j,Y^a=1).
\end{aligned}
\]
Summing the analogous identity over both values of \(Y^a\) yields
\[
P(X=j,A=a)=\frac12P(X=j)=\frac1{2d},
\]
where the last equality uses \cref{obj:ass:uniform-covariates}.

By \cref{obj:ass:consistency}, the following indicator identity holds almost surely:
\[
\begin{aligned}
\mathbf1_{\{X=j\}}S
={}&2\mathbf1_{\{X=j,A=1,Y^1=1\}}
-\mathbf1_{\{X=j,A=1\}}\\
&-2\mathbf1_{\{X=j,A=0,Y^0=1\}}
+\mathbf1_{\{X=j,A=0\}}.
\end{aligned}
\]
Integrating and using the fair-assignment identities gives
\[
\mathbb E_P[\mathbf1_{\{X=j\}}S]
=P(X=j,Y^1=1)-P(X=j,Y^0=1).
\]
Dividing by \(P(X=j)=1/d>0\), we obtain
\[
\mathbb E_P[S\mid X=j]
=\mu_{1j}(P)-\mu_{0j}(P)
=\tau_j(P).
\]
% lean: CausalSmith.Stat.LdpOptvalueUniformFrontier.signedCellMean_eq_contrast

\item \emph{Identification of the paired law.}
Because \(S\) takes values in \(\{-1,1\}\), its two signed atoms satisfy
\[
\begin{aligned}
P(X=j,S=1)+P(X=j,S=-1)&=\frac1d,\\
P(X=j,S=1)-P(X=j,S=-1)
&=\mathbb E_P[\mathbf1_{\{X=j\}}S]
=\frac{\tau_j(P)}d.
\end{aligned}
\]
Solving these two equations gives, for \(s\in\{-1,1\}\),
\[
P(X=j,S=s)=\frac{1+s\tau_j(P)}{2d}.
\]
Moreover, \cref{obj:ass:interior-means} gives
\[
-\frac12
=\frac14-\frac34
\le \tau_j(P)
\le \frac34-\frac14
=\frac12.
\]
For any \(\theta\in[-1/2,1/2]^d\), the paired masses satisfy
\[
\frac{1+s\theta_j}{2d}\ge0,
\qquad
\sum_{j=1}^d\sum_{s\in\{-1,1\}}
\frac{1+s\theta_j}{2d}
=\sum_{j=1}^d\frac1d=1.
\]
Thus \(P_{\tau(P)}\) is a probability law. Equality of the displayed atom probabilities on the finite alphabet proves
\[
\mathcal L_P(X,S)=P_{\tau(P)}.
\]
% lean: CausalSmith.Stat.LdpOptvalueUniformFrontier.causal_signed_law

\item \emph{Decomposition of the causal value.}
For each stratum, the two cases
\(\mu_{0j}(P)\le\mu_{1j}(P)\) and
\(\mu_{1j}(P)\le\mu_{0j}(P)\) give
\[
\max\{\mu_{0j}(P),\mu_{1j}(P)\}
=\frac{\mu_{0j}(P)+\mu_{1j}(P)}2
+\frac{|\tau_j(P)|}2.
\]
Consequently,
\[
V(P)
=\frac1d\sum_{j=1}^d
 \frac{\mu_{0j}(P)+\mu_{1j}(P)}2
+\frac{F(\tau(P))}2.
\]

To identify the first term, consistency gives the almost-sure identity
\[
\mathbf1_{\{X=j\}}Y
=\mathbf1_{\{X=j,A=0,Y^0=1\}}
+\mathbf1_{\{X=j,A=1,Y^1=1\}}.
\]
The fair-assignment identity
\(P(X=j,A=a,Y^a=1)=\frac12P(X=j,Y^a=1)\)
therefore yields
\[
\begin{aligned}
\mathbb E_P[\mathbf1_{\{X=j\}}Y]
&=\frac{P(X=j,Y^0=1)+P(X=j,Y^1=1)}2\\
&=\frac1d\,\frac{\mu_{0j}(P)+\mu_{1j}(P)}2.
\end{aligned}
\]
Since the strata partition the sample space,
\[
B(P)=\mathbb E_P[Y]
=\sum_{j=1}^d\mathbb E_P[\mathbf1_{\{X=j\}}Y]
=\frac1d\sum_{j=1}^d
 \frac{\mu_{0j}(P)+\mu_{1j}(P)}2.
\]
Substitution proves
\[
V(P)=B(P)+\frac{F(\tau(P))}2.
\]
% lean: CausalSmith.Stat.LdpOptvalueUniformFrontier.causal_value_decomposition

\item \emph{Total variation from the uniform law.}
Fix \(\theta\in[-1/2,1/2]^d\). Both \(P_\theta\) and
\(\mathsf U_{2d}=P_0\) are probability laws by the mass calculation above.
Define the event of nonnegative atom differences by
\[
\mathcal E_{\theta,+}
:=\left\{(j,s)\in[d]\times\{-1,1\}:
P_\theta(j,s)-\frac1{2d}\ge0\right\}.
\]
The atom differences sum to zero:
\[
\sum_{j=1}^d\sum_{s\in\{-1,1\}}
\left(P_\theta(j,s)-\frac1{2d}\right)=0.
\]
Splitting this sum into its nonnegative and negative terms gives
\[
P_\theta(\mathcal E_{\theta,+})
-\mathsf U_{2d}(\mathcal E_{\theta,+})
=\frac12\sum_{j=1}^d\sum_{s\in\{-1,1\}}
\left|P_\theta(j,s)-\frac1{2d}\right|.
\]
For every event \(\mathcal E\subseteq[d]\times\{-1,1\}\), its sum of atom
differences lies between the sum of all negative differences and the sum
of all nonnegative differences. Hence
\[
\left|P_\theta(\mathcal E)-\mathsf U_{2d}(\mathcal E)\right|
\le\frac12\sum_{j=1}^d\sum_{s\in\{-1,1\}}
\left|P_\theta(j,s)-\frac1{2d}\right|.
\]
The event \(\mathcal E_{\theta,+}\) attains this bound, proving the finite
atom formula for total variation. For each \(j\), its two summands reduce to
\[
\begin{aligned}
\sum_{s\in\{-1,1\}}
\left|\frac{1+s\theta_j}{2d}-\frac1{2d}\right|
&=\left|\frac{\theta_j}{2d}\right|
+\left|-\frac{\theta_j}{2d}\right|\\
&=\frac{|\theta_j|}{d}.
\end{aligned}
\]
Therefore
\[
\operatorname{TV}(P_\theta,\mathsf U_{2d})
=\frac1{2d}\sum_{j=1}^d|\theta_j|
=\frac{F(\theta)}2.
\]
Applying this identity at \(\theta=\tau(P)\) proves
\[
V(P)=B(P)+\operatorname{TV}(P_{\tau(P)},\mathsf U_{2d}).
\]
% lean: CausalSmith.Stat.LdpOptvalueUniformFrontier.tvFromUniform_pairedLaw

\item \emph{The symmetric causal construction.}
Fix \(\theta\in[-1/2,1/2]^d\). For
\(a,y\in\{0,1\}\) and \(j\in[d]\), define the Bernoulli masses
\[
\pi_{aj}(y)
:=
\left(\frac{1+(2a-1)\theta_j}{2}\right)^y
\left(\frac{1-(2a-1)\theta_j}{2}\right)^{1-y}.
\]
Their success probabilities belong to \([1/4,3/4]\), so
\[
\pi_{aj}(y)\ge0,
\qquad
\sum_{y\in\{0,1\}}\pi_{aj}(y)=1.
\]
By \cref{obj:def:symmetric-law}, the full-data atom masses are
\[
\begin{aligned}
&G_\theta(X=j,A=a,Y=y,Y^0=y^0,Y^1=y^1)\\
&\qquad=\frac1{2d}\,
\pi_{0j}(y^0)\pi_{1j}(y^1)
\mathbf1_{\{y=(1-a)y^0+ay^1\}},
\qquad a,y,y^0,y^1\in\{0,1\}.
\end{aligned}
\]
Summing first over \(y\), then over the two potential outcomes and
assignment, gives total mass
\[
\sum_{j=1}^d\sum_{a\in\{0,1\}}\frac1{2d}
\left(\sum_{y^0\in\{0,1\}}\pi_{0j}(y^0)\right)
\left(\sum_{y^1\in\{0,1\}}\pi_{1j}(y^1)\right)
=1.
\]
Thus \(G_\theta\) is a probability law. The same atom formula gives
\[
G_\theta(X=j)=\frac1d
\]
and
\[
\begin{aligned}
&G_\theta(X=j,A=a,Y^0=y^0,Y^1=y^1)\\
&\qquad=\frac1{2d}\pi_{0j}(y^0)\pi_{1j}(y^1)
=\frac12G_\theta(X=j,Y^0=y^0,Y^1=y^1).
\end{aligned}
\]
These are the conditions in
\cref{obj:ass:uniform-covariates,obj:ass:fair-randomization}.
The atom formula assigns all mass to \(Y=Y_A^{\mathrm{pot}}\), supplying
\cref{obj:ass:consistency}. Finally, summing over the other coordinates yields
\[
G_\theta(X=j,Y^a=1)
=\frac1d\,\frac{1+(2a-1)\theta_j}{2},
\]
and hence
\[
\mu_{0j}(G_\theta)=\frac{1-\theta_j}{2},
\qquad
\mu_{1j}(G_\theta)=\frac{1+\theta_j}{2}.
\]
Both means lie in \([1/4,3/4]\), as required by
\cref{obj:ass:interior-means}. All four conditions of
\cref{obj:def:causal-model} therefore hold.
% lean: CausalSmith.Stat.LdpOptvalueUniformFrontier.symmetricLaw_causalModel

\item \emph{Symmetric value and the ancillary assignment.}
The arm means just computed satisfy
\[
\tau_j(G_\theta)=\theta_j,
\qquad
\frac{\mu_{0j}(G_\theta)+\mu_{1j}(G_\theta)}2=\frac12.
\]
Applying the armwise maximum decomposition gives
\[
V(G_\theta)
=\frac1d\sum_{j=1}^d\frac12+\frac{F(\theta)}2
=\frac12+\frac{F(\theta)}2.
\]
% lean: CausalSmith.Stat.LdpOptvalueUniformFrontier.symmetricLaw_value

Define the observed-record reconstruction map on the whole paired
alphabet and both assignment values by
\[
\operatorname{rec}(j,s,a)
:=\left(j,a,\frac{1+s(2a-1)}2\right),
\qquad
j\in[d],\quad s\in\{-1,1\},\quad a\in\{0,1\}.
\]
Its outcome coordinate is binary. Summing the full-data atom masses over
the unused potential outcome gives
\[
\begin{aligned}
G_\theta(X=j,S=s,A=a)
&=\frac1{2d}\,
\pi_{aj}\!\left(\frac{1+s(2a-1)}2\right)\\
&=\frac1{2d}\,\frac{1+s\theta_j}{2}
=\frac{P_\theta(j,s)}2.
\end{aligned}
\]
Indeed, for \(a=0\), the reconstructed outcome is \((1-s)/2\), so the
two choices of \(s\) select the failure and success masses of
\(\operatorname{Bernoulli}((1-\theta_j)/2)\). For \(a=1\), it is
\((1+s)/2\), selecting the success and failure masses of
\(\operatorname{Bernoulli}((1+\theta_j)/2)\). In both cases the selected
mass is \((1+s\theta_j)/2\). Thus the joint signed-input and assignment
law is
\[
\mathcal L_{G_\theta}(X,S,A)
=P_\theta\otimes\operatorname{Bernoulli}(1/2).
\]
% lean: CausalSmith.Stat.LdpOptvalueUniformFrontier.symmetricLaw_ancillary

For every binary observed record, \((2A-1)^2=1\), and therefore
\[
S(2A-1)
=(2A-1)(2Y-1)(2A-1)=2Y-1.
\]
This identity holds pointwise and hence \(G_\theta\)-almost surely.
% lean: CausalSmith.Stat.LdpOptvalueUniformFrontier.outcome_sign_reconstruction

\item \emph{Class preservation by averaging.}
Fix \(\varepsilon\in\mathbb R\) and an original-record protocol \(Q\).
Using the same seed law and message spaces, define
\[
\mathsf K_i(\cdot\mid(j,s),h,r)
:=\frac12\sum_{a\in\{0,1\}}
Q_i(\cdot\mid\operatorname{rec}(j,s,a),h,r).
\]
Each kernel is a finite mixture of measurable probability kernels, so
this defines a protocol.

Suppose \(Q\) satisfies \cref{obj:ass:full-record-privacy}. For every
\(i,h,r\), measurable message event \(E\), and paired inputs
\((j,s),(j',s')\), apply that inequality separately with assignment zero
and assignment one:
\[
Q_i(E\mid\operatorname{rec}(j,s,a),h,r)
\le e^\varepsilon
Q_i(E\mid\operatorname{rec}(j',s',a),h,r),
\qquad a\in\{0,1\}.
\]
Adding the two inequalities with weights \(1/2\) gives
\[
\mathsf K_i(E\mid(j,s),h,r)
\le e^\varepsilon\mathsf K_i(E\mid(j',s'),h,r).
\]
This proves sequential membership as specified by
\cref{obj:def:sequential-class}. If \(Q\) also satisfies
\cref{obj:ass:noninteraction}, then for all histories \(h,h'\),
\[
\begin{aligned}
\mathsf K_i(\cdot\mid(j,s),h,r)
&=\frac12\sum_a
Q_i(\cdot\mid\operatorname{rec}(j,s,a),h,r)\\
&=\frac12\sum_a
Q_i(\cdot\mid\operatorname{rec}(j,s,a),h',r)
=\mathsf K_i(\cdot\mid(j,s),h',r).
\end{aligned}
\]
Thus \cref{obj:def:noninteractive-class} is preserved as well. These
arguments apply at every declared history and seed, for every real
\(\varepsilon\).
% lean: CausalSmith.Stat.LdpOptvalueUniformFrontier.signed_protocol_classes

\item \emph{Reconstruction of the complete sampled input.}
For deterministic vectors, define
\[
\begin{gathered}
\mathbf v=((j_i,s_i))_{i=1}^n
\in([d]\times\{-1,1\})^n,
\qquad
\mathbf a=(a_i)_{i=1}^n\in\{0,1\}^n,\\
\operatorname{rec}^{(n)}(\mathbf v,\mathbf a)
:=\bigl(\operatorname{rec}(j_i,s_i,a_i)\bigr)_{i=1}^n.
\end{gathered}
\]
The single-participant ancillary identity and
\cref{obj:ass:iid-people} imply
\[
\mathcal L_{G_\theta}\!\left(
((X_i,S_i))_{i=1}^n,(A_i)_{i=1}^n\right)
=P_\theta^{\otimes n}
 \otimes\operatorname{Bernoulli}(1/2)^{\otimes n}.
\]
To see the factorization directly, the probability of a specified pair
of vectors is
\[
\prod_{i=1}^n\frac{P_\theta(j_i,s_i)}2
=\left(\prod_{i=1}^nP_\theta(j_i,s_i)\right)2^{-n}.
\]
Reconstruction recovers every observed record, so pushing this product
law through \(\operatorname{rec}^{(n)}\) gives
\[
\bigl(P_\theta^{\otimes n}
 \otimes\operatorname{Bernoulli}(1/2)^{\otimes n}\bigr)
 \circ(\operatorname{rec}^{(n)})^{-1}
=(G_\theta)_O^{\otimes n}.
\]

Write \(\lambda=\mathcal L(R)\) for the protocol's seed law. Including
the independent randomness from
\cref{obj:ass:independent-randomness}, define
\[
\Lambda_\theta
:=P_\theta^{\otimes n}
 \otimes\operatorname{Bernoulli}(1/2)^{\otimes n}
 \otimes\lambda\otimes\operatorname{Uniform}[0,1]
\]
on tuples \((\mathbf v,\mathbf a,r,u)\). Its pushforward satisfies
\[
\mathcal L_{G_\theta}(\mathbf O,R,U)
=\Lambda_\theta\circ
\bigl[(\mathbf v,\mathbf a,r,u)
\mapsto(\operatorname{rec}^{(n)}(\mathbf v,\mathbf a),r,u)\bigr]^{-1}.
\]
Thus running \(Q\) on the sampled causal records has the same joint
message and randomizer law as running it on these reconstructed
product-law inputs.
% lean: CausalSmith.Stat.LdpOptvalueUniformFrontier.symmetricLaw_reconstructed_decisionLaw

\item \emph{Averaging through the sequential recursion.}
For \(0\le\ell\le n\), define the prefix message space
\[
\mathcal Z_{\le\ell}:=\prod_{i=1}^{\ell}\mathcal Z_i,
\qquad
\mathcal Z_{\le0}:=\{\langle\rangle\}.
\]
For a fixed observed-input vector
\(\mathbf o\in([d]\times\{0,1\}^2)^n\), paired-input vector
\(\mathbf v\in([d]\times\{-1,1\})^n\), and seed \(r\in\mathcal R\),
let \(\Pi_\ell^Q(\cdot\mid\mathbf o,r)\) and
\(\Pi_\ell^{\mathsf K}(\cdot\mid\mathbf v,r)\) be the prefix probability
laws defined by
\[
\Pi_0^Q=\Pi_0^{\mathsf K}=\delta_{\langle\rangle}
\]
and, for \(0\le\ell<n\),
\[
\begin{aligned}
\Pi_{\ell+1}^Q(dh,dz\mid\mathbf o,r)
&=\Pi_\ell^Q(dh\mid\mathbf o,r)\,
 Q_{\ell+1}(dz\mid o_{\ell+1},h,r),\\
\Pi_{\ell+1}^{\mathsf K}(dh,dz\mid\mathbf v,r)
&=\Pi_\ell^{\mathsf K}(dh\mid\mathbf v,r)\,
 \mathsf K_{\ell+1}(dz\mid v_{\ell+1},h,r),
\end{aligned}
\]
where \(h\in\mathcal Z_{\le\ell}\), \(z\in\mathcal Z_{\ell+1}\), and
\((h,z)\) denotes the appended history.

We prove by induction on \(\ell\) that, for every measurable
\(\varphi:\mathcal Z_{\le\ell}\to[0,\infty]\),
\[
\begin{aligned}
&2^{-n}\sum_{\mathbf a\in\{0,1\}^n}
 \int\varphi(h)\,
 \Pi_\ell^Q(dh\mid
 \operatorname{rec}^{(n)}(\mathbf v,\mathbf a),r)\\
&\qquad=\int\varphi(h)\,
 \Pi_\ell^{\mathsf K}(dh\mid\mathbf v,r).
\end{aligned}
\]
For \(\ell=0\), both prefix laws are concentrated on the empty history,
and \(2^{-n}\sum_{\mathbf a\in\{0,1\}^n}1=1\).

For the induction step, take measurable
\(\varphi:\mathcal Z_{\le\ell+1}\to[0,\infty]\) and define
\[
\varphi_{\mathrm{av},\ell}(h)
:=\int\varphi(h,z)\,
 \mathsf K_{\ell+1}(dz\mid v_{\ell+1},h,r),
\qquad h\in\mathcal Z_{\le\ell}.
\]
This is a measurable nonnegative function. For \(i\in[n]\) and
\(a\in\{0,1\}\), define the updated bit vector by
\[
(\mathbf a^{\,i\leftarrow a})_{i'}
:=
\begin{cases}
a,&i'=i,\\
a_{i'},&i'\ne i,
\end{cases}
\qquad i'\in[n].
\]
The uniform distribution on bit vectors is invariant under flipping
coordinate \(i\). Pairing each vector with its flipped vector shows that
averaging any nonnegative function over all bit vectors equals
averaging its two versions with coordinate \(i\) set to zero and one.

Apply this identity at \(i=\ell+1\). Updating that coordinate leaves the
first \(\ell\) inputs unchanged, hence leaves the preceding prefix law
unchanged. The recursion and the definition of \(\mathsf K_{\ell+1}\)
therefore give, for every \(\mathbf a\),
\[
\begin{aligned}
&\frac12\sum_{a\in\{0,1\}}
 \int\varphi(h,z)\,
 \Pi_{\ell+1}^Q
 \bigl(dh,dz\mid
 \operatorname{rec}^{(n)}
   (\mathbf v,\mathbf a^{\,\ell+1\leftarrow a}),r\bigr)\\
&\qquad=
 \int\left[
 \frac12\sum_{a\in\{0,1\}}
 \int\varphi(h,z)\,
 Q_{\ell+1}(dz\mid
 \operatorname{rec}(j_{\ell+1},s_{\ell+1},a),h,r)
 \right]\\
&\hspace{6em}\cdot
 \Pi_\ell^Q(dh\mid
 \operatorname{rec}^{(n)}(\mathbf v,\mathbf a),r)\\
&\qquad=
 \int\varphi_{\mathrm{av},\ell}(h)\,
 \Pi_\ell^Q(dh\mid
 \operatorname{rec}^{(n)}(\mathbf v,\mathbf a),r).
\end{aligned}
\]
Averaging this equality over \(\mathbf a\), applying the induction
hypothesis to \(\varphi_{\mathrm{av},\ell}\), and using the paired
recursion proves
\[
\begin{aligned}
&2^{-n}\sum_{\mathbf a}
 \int\varphi\,
 \Pi_{\ell+1}^Q
 (\,\cdot\mid\operatorname{rec}^{(n)}(\mathbf v,\mathbf a),r)\\
&\qquad=
 \int\varphi_{\mathrm{av},\ell}(h)\,
 \Pi_\ell^{\mathsf K}(dh\mid\mathbf v,r)
 =\int\varphi\,\Pi_{\ell+1}^{\mathsf K}
 (\,\cdot\mid\mathbf v,r).
\end{aligned}
\]
At \(\ell=n\), indicator functions of measurable transcript events give
the equality of probability laws
\[
2^{-n}\sum_{\mathbf a\in\{0,1\}^n}
\Pi_n^Q
(\,\cdot\mid\operatorname{rec}^{(n)}(\mathbf v,\mathbf a),r)
=\Pi_n^{\mathsf K}(\,\cdot\mid\mathbf v,r).
\]
Integrating against \(P_\theta^{\otimes n}\), \(\lambda\), and the uniform
analyst randomizer, and using the identity expressing
\(\mathcal L_{G_\theta}(\mathbf O,R,U)\) as the reconstruction
pushforward of \(\Lambda_\theta\), proves
\[
\begin{aligned}
&\mathcal L_{G_\theta,Q}(\mathbf Z,R,U)\\
&\qquad=
P_\theta^{\otimes n}\otimes\lambda
 \otimes\operatorname{Uniform}[0,1]
\ \text{followed by }\ 
\mathbf Z\sim\Pi_n^{\mathsf K}(\cdot\mid\mathbf v,r).
\end{aligned}
\]
Equivalently, for every measurable
\(\mathcal E\subseteq\mathcal Z_{\le n}\times\mathcal R\times[0,1]\),
the common probability is
\[
\int\!\int\!\int\!\int
\mathbf1_{\mathcal E}(\mathbf z,r,u)\,
\Pi_n^{\mathsf K}(d\mathbf z\mid\mathbf v,r)\,
\operatorname{Uniform}[0,1](du)\,
\lambda(dr)\,P_\theta^{\otimes n}(d\mathbf v).
\]
Projecting onto \((R,\mathbf Z)\) gives the first asserted joint
seed-transcript equality. The induction begins at zero and therefore
also covers \(n=0\).
% lean: CausalSmith.Stat.LdpOptvalueUniformFrontier.averagedProtocol_reconstructed_decisionLaw

\item \emph{Class preservation by signed preprocessing.}
For the converse direction, fix a paired-alphabet protocol \(\mathsf K\).
Define the deterministic paired-record map by
\[
\operatorname{pair}(j,a,y)
:=\bigl(j,(2a-1)(2y-1)\bigr),
\qquad
(j,a,y)\in[d]\times\{0,1\}^2,
\]
and define \(Q\), with the same seed law and message spaces, by
\[
Q_i(\cdot\mid(j,a,y),h,r)
:=\mathsf K_i(\cdot\mid\operatorname{pair}(j,a,y),h,r).
\]
Composition with this measurable finite-alphabet map preserves the
probability-kernel property.

If \(\mathsf K\) satisfies the paired privacy inequality, then for every
two observed records \(o,o'\), every history and seed, and every
measurable message event \(E\),
\[
\begin{aligned}
Q_i(E\mid o,h,r)
&=\mathsf K_i(E\mid\operatorname{pair}(o),h,r)\\
&\le e^\varepsilon
 \mathsf K_i(E\mid\operatorname{pair}(o'),h,r)
=e^\varepsilon Q_i(E\mid o',h,r).
\end{aligned}
\]
Thus \(Q\) satisfies \cref{obj:ass:full-record-privacy} and belongs to
the sequential class whenever \(\mathsf K\) does. If \(\mathsf K\) is
independent of history, then
\[
Q_i(\cdot\mid o,h,r)
=\mathsf K_i(\cdot\mid\operatorname{pair}(o),h,r)
=\mathsf K_i(\cdot\mid\operatorname{pair}(o),h',r)
=Q_i(\cdot\mid o,h',r).
\]
This supplies \cref{obj:ass:noninteraction} and hence the noninteractive
membership in \cref{obj:def:noninteractive-class}.
% lean: CausalSmith.Stat.LdpOptvalueUniformFrontier.signed_protocol_classes

\item \emph{The converse transcript identity.}
For a deterministic observed vector \(\mathbf o\), define its
coordinatewise paired image by
\[
\operatorname{pair}^{(n)}(\mathbf o)
:=\bigl(\operatorname{pair}(o_i)\bigr)_{i=1}^n.
\]
The previously established identities
\(\tau(G_\theta)=\theta\) and
\(\mathcal L_{G_\theta}(X,S)=P_\theta\) imply, by coordinatewise
pushforward of the independent observed sample,
\[
\begin{aligned}
&\bigl((G_\theta)_O^{\otimes n}
 \otimes\lambda\otimes\operatorname{Uniform}[0,1]\bigr)
 \circ
 \bigl[(\mathbf o,r,u)
 \mapsto(\operatorname{pair}^{(n)}(\mathbf o),r,u)\bigr]^{-1}\\
&\qquad=
P_\theta^{\otimes n}\otimes\lambda
 \otimes\operatorname{Uniform}[0,1].
\end{aligned}
\]
Here \cref{obj:ass:iid-people,obj:ass:independent-randomness} identify the
measure on the left before pushforward with the sampling scheme's
joint input law.

For each fixed \(\mathbf o,r\), induction through the prefix recursion
gives
\[
\Pi_\ell^Q(\cdot\mid\mathbf o,r)
=\Pi_\ell^{\mathsf K}
(\cdot\mid\operatorname{pair}^{(n)}(\mathbf o),r),
\qquad 0\le\ell\le n.
\]
The zero-length laws are the same point mass. At a successor step,
the induction hypothesis identifies the preceding prefix laws, and
the defining equality
\[
Q_{\ell+1}(dz\mid o_{\ell+1},h,r)
=\mathsf K_{\ell+1}
(dz\mid\operatorname{pair}(o_{\ell+1}),h,r)
\]
identifies their extension kernels. Thus their successor prefix laws
are equal as well. In particular, their full fixed-input transcript
laws agree at \(\ell=n\).

Integrating this fixed-input equality and using the displayed input
pushforward proves that the joint law of \((\mathbf Z,R,U)\) under
\(G_\theta,Q\) equals the law obtained from independent \(P_\theta\)
inputs, the seed law \(\lambda\), and \(\mathsf K\). Projecting onto
\((R,\mathbf Z)\) gives the asserted converse equality. The same
zero-length base case covers \(n=0\).
% lean: CausalSmith.Stat.LdpOptvalueUniformFrontier.pulledProtocol_decisionLaw
\end{enumerate}
\end{proof}
